% Held out of the report (2026-11-26, user): older DRTC results wait for the new robot results (E4).
% Not built; kept unchanged for later.

\section{Why the shipped client fails under faults}

The shipped client has two give-up limits (Section~\ref{sec:baseline}). The starvation grace ends a run one second
after the queue runs dry, and every fault of the mix lasts about that long or longer; worse, when the queue runs dry
the request in flight is forgotten, so the late reply that would have saved the run is discarded. In the first run
every faulted trial ended this way at the first fault (Section~\ref{sec:main}), and on the bench each fault alone was
enough (Section~\ref{sec:bench}). The failures therefore came from the client's rules, not from the policies, which
succeeded in most trials on the normal network.

Raising the limits removes the give-up but not the weakness behind it. The client still keeps only half a second of
actions in reserve, and it replaces a lost request only once, when its queue runs dry: if that fresh request is lost
too, nothing more is sent, and the arm waits until LeRobot's 10~s startup deadline ends the run. In our local replay checks a patient client stalled for more than half of a faulted run, with a single
stall of 7.5~s after the dropped replies. This is why setup~v2 compares \drtc{} with the patient client, and why its
primary comparisons use the stall fraction and the recovery time rather than whether the run was ended.

\section{Why \drtc{} recovers}

\drtc{} changes three things that matter under faults. It asks earlier (with $H-s_{\min}$ actions left), so a short
delay is absorbed by the queue; it holds the arm when the queue is empty instead of ending the run; and its cooldown
replaces an overdue reply with a fresh request, so the first chunk after a fault is computed from the current scene.
Because the policies are closed-loop, they continue the task from wherever the arm stopped. In the first run this
brought success back to a level not significantly different from the fault-free baseline for both policies, and the
arm moved again within 0--1.5~s of every fault's end (Table~\ref{tab:mainrecovery}). Setup~v2 separated these parts
(Section~\ref{sec:v2bench}): against a client that also waits, \drtc's advantage comes almost entirely from the
cooldown re-request, and the same rule makes it worse when the link only slows down, because it gives up replies that
are late but coming.

\section{Short-chunk policies}\label{sec:shortchunk}

\groot's chunk holds 16 actions, 0.53~s at 30~Hz, barely longer than its round trip of about 0.45~s. Even on the
normal network its queue empties before the next chunk is ready, and the arm stood still in
\R{main/groot/baseline/M2} of the ticks. This is a property of the combination of chunk length and latency rather than
of the network faults, and it is where adaptive playback helps most: slowing down instead of stopping cut the stall
fraction to \R{main/groot/smooth/M2}. The cost is speed: slower playback lengthens the trial a little (time to success
\R{main/groot/drtc/M1t} with \drtc{} alone, \R{main/groot/smooth/M1t} with smoothing). For policies with long chunks,
such as \smolvla{} with 50 actions, the queue rarely runs dry and playback stays at full rate. The same arithmetic
explains our first fix: a cooldown capped at half a chunk gives up every request when the round trip is longer than
half a chunk, which for a 16-action policy over WiFi is always.

\section{Smoothness and success}

In the first run the filter and adaptive playback made the motion markedly smoother (jerk lower by factors of
\R{main/smolvla/jerkratio} and \R{main/groot/jerkratio}) with no loss of success; the success rates with smoothing were
higher than with \drtc{} alone for both policies, but with ten trials per condition that difference is not
significant. Two cautions apply. First, a low jerk can come from standing still: the first version of the metric
counted the time the arm held still, which made the stalling client look smooth; metrics version~3 counts only the
moving time. Second, smoothing can lag behind the policy: a causal filter delays the motion, and the bench has no task
success to show the cost. Setup~v2 therefore measures the deviation from the policy's plan (M12) and refuses a smoothing
method that adds more than one degree of it.

\section{Smoothing on a slow link}\label{sec:smoothdisc}

Setup~v2 separated the smoothing methods (Sections~\ref{sec:v2bench} and~\ref{sec:v2e3b}), and three lessons stand out.

\paragraph{Match the method to the source.} Each method removed the source it was designed for and little else: the
Butterworth filter the jitter and the seams (most of the jerk), adaptive playback the stalls of the policies whose
chunk is barely longer than their round trip. Which policies need what is decided by the ratio of chunk length to
round trip, not by the model's size: ACT (100 actions) and \smolvla{} (50) stalled for 3\% and 12--15\% of a faulted
run, the two GR00T models (16 actions) for about two thirds of it without adaptive playback.

\paragraph{RTC is for inference-bound links.} Guided real-time chunking makes the seam smooth at the source, but it
asks the server for more work per request. When the round trip is mostly network (about 260~ms of our 300--450~ms),
the extra time per request costs more stalls than the smoother seam saves: \smolvla's stalls doubled. RTC was chosen
by the pre-registered rule on a single policy-task pair, and the ablation of E3b shows that the stack is better
without it. We keep the rule as fixed and report the ablation as its correction; on a link where inference dominates
the round trip, the trade-off may well reverse.

\paragraph{Smoothness has a price, and it has to be measured.} The full stack moves the arm 0.7 to 2.3\degree{} away
from the policy's own plan on average. On the bench this cannot cost a task, because the replay robot does not grasp;
on the robot it might. This is why the deviation (M12) is measured and limited in the stack rule, and why the robot
experiment E4 is the test that matters for the stack: it is the only one with task success.

